% Generated draft: requires mathematical and source review.
\documentclass[11pt]{article}
\usepackage[margin=25mm]{geometry}
\usepackage{fontspec}
\setmainfont{Latin Modern Roman}
\usepackage{amsmath,amssymb,mathtools,unicode-math}
\setmathfont{Latin Modern Math}
\usepackage{enumitem,needspace,xurl,hyperref}
\setlength{\parindent}{0pt}
\setlength{\parskip}{5pt}
\setlength{\emergencystretch}{3em}
\newcommand{\N}{\mathbb{N}}
\newcommand{\Z}{\mathbb{Z}}
\newcommand{\Q}{\mathbb{Q}}
\newcommand{\R}{\mathbb{R}}
\newcommand{\C}{\mathbb{C}}
\newcommand{\F}{\mathbb{F}}
\newcommand{\A}{\mathbb{A}}
\newcommand{\PP}{\mathbb{P}}
\newcommand{\E}{\mathbb{E}}
\newcommand{\eps}{\varepsilon}
\newcommand{\dd}{\,\mathrm{d}}
\begin{document}
\section*{UnsolvedMath ID: 800198}
\textbf{Source problem code:} catalogue entry
\subsection{Definitions and mathematical statement}
The supplied question does not specify its conventions for hyperk\"ahler manifolds, deformations, or Lagrangian fibrations. The following is a provisional formulation under the compact irreducible holomorphic symplectic convention; these conventions require confirmation against the source.

A hyperk\"ahler manifold in this convention is a compact, simply connected K\"ahler manifold \(X\) of complex dimension \(2n\), where \(n\geq 1\), such that
\[
H^0(X,\Omega_X^2)=\C\sigma_X
\]
for an everywhere nondegenerate holomorphic two-form \(\sigma_X\). Two such manifolds are deformation equivalent if they can be connected by a finite chain of smooth proper holomorphic families over connected complex bases, with successive manifolds occurring as fibers of a common family.

For this formulation, a holomorphic Lagrangian fibration on such a manifold \(Y\) is a surjective holomorphic map \(f:Y\to B\) with connected fibers, where \(B\) is a normal compact complex space of dimension \(n\), for which there is a dense open subset \(B^\circ\subseteq B\) over which \(f\) is smooth and every fiber \(F\) satisfies
\[
\dim_{\C}F=n,\qquad \sigma_Y|_F=0.
\]

\textbf{Question (Beauville, 2010, Question 7).} Is it true that for every hyperk\"ahler manifold \(X\), there exists a deformation-equivalent hyperk\"ahler manifold \(Y\) admitting a holomorphic Lagrangian fibration?

The compact irreducible interpretation, the deformation-equivalence convention, and the stated conventions concerning the base and fibers are not explicit in the supplied record. Thus this precise formulation remains conditional on those interpretations. The record supplies no projectivity hypothesis on \(X\) or \(Y\), and no prescribed isomorphism type for \(B\).
\subsection{Short English statement}
Beauville asks whether every hyperk\"ahler manifold can be deformed to one admitting a Lagrangian fibration. Under the provisional compact irreducible interpretation, the general fibers have half the complex dimension of the manifold and the holomorphic symplectic form restricts to zero on them. The source record leaves several conventions unspecified.
\subsection{Sources}
\begin{itemize}
\item Local UnsolvedMath registry, permanent ID \(800198\), source problem code \texttt{catalogue entry}, titled ``Holomorphic Symplectic Geometry --- Question 7.'' The supplied record attributes the question to Beauville (2010) and labels it \texttt{open}; this status is supplied metadata, not an independently verified mathematical or literature assessment.
\item Beauville, \emph{Holomorphic Symplectic Geometry: A Problem List} (2010), Question 7. Supplied source URL: \url{https://arxiv.org/abs/1002.4321}. The background records access on 2026-10-16 and extraction by \texttt{arxiv\_latex}; these are supplied provenance details, not a claim of independent access or verification.
\end{itemize}
\end{document}
